\section{Revised Two-Architecture Overview}
\label{sec:two-arch}

10BASE-T1S is Alfred's initial wedge, but what it offers differs by customer starting point.

\begin{center}
\small
\begin{tabularx}{\linewidth}{@{}p{3.4cm} X@{}}
\toprule
\textbf{Customer} & \textbf{Value proposition} \\
\midrule
Clean-sheet & Adopt Alfred's Ethernet-native electrical infrastructure directly. Peripherals attach through MCU-less endpoint silicon (LAN8660/LAN8661) -- no per-node firmware to write, own, or update. \\
Existing vehicle & Keep the existing E/E architecture exactly as it is. Attach Alfred through Gateway hardware that observes, timestamps, and represents that architecture on an Alfred-facing 10BASE-T1S pseudo-bus, with an explicitly gated path to authorized bidirectional control. \\
\bottomrule
\end{tabularx}
\end{center}

\begin{diagram}
                        ALFRED

                    10BASE-T1S
                         |
             +-----------+-----------+
             |                       |
             v                       v

      CLEAN-SHEET                 BROWNFIELD
      NEW PROGRAM              EXISTING PROGRAM
             |                       |
       direct E2B                 Gateway
       endpoints                    |
             |                 CAN/CAN-FD/LIN
             v                       |
       peripherals               existing ECUs
\end{diagram}

\subsection{What each side is actually built from}

\begin{diagram}
   GREENFIELD SIDE                       BROWNFIELD SIDE

   Clicker 4 (temporary central node)    Clicker 4 (temporary central node)
        |                                     |
     LAN8651 (host MAC-PHY)                LAN8651 (host MAC-PHY)
        |                                     |
   10BASE-T1S ---------------------------- 10BASE-T1S
        |                                     |
   +----+----+                          Alfred Gateway
   |         |                          (MCU + LAN8651)
   v         v                                |
 LAN8660  LAN8661                    CAN / CAN-FD / LIN transceivers
   |         |                                |
 motor    LED                          existing vehicle network
 driver   driver                       (unmodified, still operating)
\end{diagram}

Both sides put a device on the \emph{same} 10BASE-T1S electrical layer, using the \emph{same} host-side MAC-PHY part (LAN8651) wherever a programmable MCU needs network access. What differs is what sits on the peripheral side of that network: fixed-function, MCU-less endpoint silicon for a clean-sheet peripheral (Section~\ref{sec:hardware-roles}), versus a programmable Gateway MCU translating an entire pre-existing legacy network (Section~\ref{sec:arch-b}). The commercial argument to a customer is exactly this symmetry: whichever side they start on, they land on the same wire, and the same Clicker~4 (standing in for Alfred's eventual central compute) sees both.
